\documentclass[12pt,a4paper]{article}

\usepackage{fontawesome5}
\providecommand{\faPencilAlt}{\faIcon{pencil-alt}}

\usepackage{styles/couleurs}
\usepackage{styles/mise_en_page}
\usepackage{styles/boites}
\usepackage{styles/maths}
\usepackage{styles/figures}

\usepackage{enumitem}
\usepackage{array}
\usepackage{tabularx}
\usepackage{amsmath}
\usepackage{amssymb}
\usepackage{tikz}
\usepackage{xparse}
\usepackage{graphicx}
\usepackage{pdflscape}
\usepackage{colortbl}

\usetikzlibrary{calc,arrows.meta,positioning,patterns}

\setlist[itemize]{leftmargin=1.4em,itemsep=0.15em,topsep=0.2em}
\setlist[enumerate]{leftmargin=*,itemsep=0.22em,topsep=0.25em}

\newcommand{\lignepointillee}{%
  \noindent\dotfill\par
  \vspace{0.45em}
}

\newcommand{\lignesreponse}[1]{%
  \begin{tcolorbox}[
    enhanced,
    colback=white,
    colframe=NeutreC!40,
    boxrule=0.5pt,
    arc=1.5mm,
    left=3mm,
    right=3mm,
    top=2mm,
    bottom=1mm,
    before skip=2pt,
    after skip=2pt
  ]
    \foreach \i in {1,...,#1}{\lignepointillee}
  \end{tcolorbox}
}

\newcommand{\imageactivite}[3]{%
  \IfFileExists{#3}{%
    \includegraphics[width=#1,height=#2,keepaspectratio]{#3}%
  }{%
    \fbox{%
      \parbox[c][#2][c]{#1}{%
        \centering\sffamily\small
        Image à intégrer\par
        \vspace{0.3em}
        \texttt{\detokenize{#3}}
      }%
    }%
  }%
}

\newcommand{\titreseq}[1]{%
\begin{center}
  {\sffamily\LARGE\bfseries\color{TitreC}
  Séquence 1 -- Statistiques}\\[0.2em]
  {\sffamily\large\bfseries\color{SavoirC}#1}
\end{center}
}

\begin{document}

\titreseq{Automatismes -- Séries courtes}

\remarque{Chaque série est prévue pour environ cinq minutes.
Elle peut être projetée ou distribuée sous forme de bande.}

\section*{Série 1 -- Lire des données brutes}

On considère la série :

\[
2\quad3\quad1\quad2\quad4\quad2\quad0\quad3\quad2\quad1.
\]

\begin{enumerate}
  \item Quel est l'effectif total ?
  \item Quel est l'effectif de la valeur \(2\) ?
  \item Quelle est la valeur la plus fréquente ?
\end{enumerate}

\section*{Série 2 -- Fréquences}

Dans une classe de \(25\) élèves, \(10\) viennent à pied.

\begin{enumerate}
  \item Écris la fréquence sous forme d'une fraction.
  \item Donne son écriture décimale.
  \item Donne le pourcentage correspondant.
\end{enumerate}

\section*{Série 3 -- Contrôler un tableau}

Un tableau d'effectifs contient :

\[
4\qquad7\qquad5\qquad3\qquad6.
\]

\begin{enumerate}
  \item Calcule l'effectif total.
  \item Les fréquences suivantes peuvent-elles convenir :
        \(16\%\), \(28\%\), \(20\%\), \(12\%\), \(24\%\) ?
\end{enumerate}

\section*{Série 4 -- Moyenne simple}

Calcule puis interprète mentalement la moyenne de :

\[
8\qquad10\qquad12\qquad14.
\]

\section*{Série 5 -- Retrouver une valeur}

La moyenne de cinq valeurs est \(12\).
Quatre de ces valeurs sont :

\[
9\qquad11\qquad13\qquad15.
\]

Quelle est la cinquième valeur ?

\section*{Série 6 -- Moyenne avec effectifs}

\[
\begin{array}{c|ccc}
\text{Valeur} & 5 & 10 & 15\\
\hline
\text{Effectif} & 2 & 3 & 1
\end{array}
\]

Calcule la moyenne.

\section*{Série 7 -- Médiane, effectif impair}

Détermine puis interprète la médiane :

\[
9\qquad4\qquad12\qquad7\qquad10.
\]

\section*{Série 8 -- Médiane, effectif pair}

Détermine la médiane :

\[
3\qquad8\qquad5\qquad11\qquad6\qquad14.
\]

\section*{Série 9 -- Médiane à partir d'effectifs}

\[
\begin{array}{c|ccccc}
\text{Valeur} & 1 & 2 & 3 & 4 & 5\\
\hline
\text{Effectif} & 2 & 4 & 5 & 3 & 1
\end{array}
\]

\begin{enumerate}
  \item Quel est l'effectif total ?
  \item Quelle position occupe la médiane ?
  \item Détermine la médiane.
\end{enumerate}

\section*{Série 10 -- Pourcentages et angles}

Complète mentalement.

\[
50\%\longleftrightarrow\ldots^\circ
\qquad
25\%\longleftrightarrow\ldots^\circ
\]

\[
10\%\longleftrightarrow\ldots^\circ
\qquad
20\%\longleftrightarrow\ldots^\circ
\]

\section*{Série 11 -- Angle d'un secteur}

Dans une enquête portant sur \(40\) personnes,
\(14\) choisissent la réponse A.

Calcule l'angle du secteur correspondant.

\section*{Série 12 -- Moyenne ou médiane ?}

Une entreprise a les salaires suivants :

\[
1500\quad1550\quad1600\quad1650\quad1700\quad9000.
\]

Sans tout calculer, quel indicateur sera le plus influencé
par le salaire de \(9000\) euros ?

\section*{Série 13 -- Vrai ou faux ?}

Justifie en une phrase.

\begin{enumerate}
  \item La somme des fréquences d'une série vaut \(100\).
  \item La médiane appartient toujours à la série.
  \item Une moyenne peut ne pas être une valeur observée.
  \item Dans un diagramme circulaire, la somme des angles vaut \(360^\circ\).
\end{enumerate}

\section*{Série 14 -- Contrôler une affirmation}

Dans une série, la médiane vaut \(17\).

Complète une phrase correcte commençant par :

\begin{quote}
« Au moins la moitié des valeurs... »
\end{quote}

\end{document}
